\section{The original auxiliary-prime supplier and residual Hecke values}
\label{sec:count-auxiliary-primes}

\status{CAP1--CAP14 prove the full original \texttt{ArPl} contract,
conditional on its unchanged \texttt{CohP} hypothesis.  This supplies
\texttt{JSP\_ACTUAL\_ARPL}, including \texttt{PhiCond}, the exact residual
subgroup, a Taylor--Wiles system for every prescribed finite list, its
actual arithmetic realization, and one character bound for all levels.
The only new published input is the unconditional Chebotarev theorem,
with its hypotheses and Frobenius convention checked below.  Neither
ordinary Hecke data nor \texttt{HasAut}, \texttt{Choice}, or the two
cohomological bounds are inferred from this construction.  CAP15--CAP17
also supply the full original \texttt{IsRes} values for any supplied
integral type, at every original \texttt{QData} level.  The type,
\texttt{TThetaEx}/\texttt{neatU}, and ordinary-control producers remain
explicit.}

Retain the original prime $\ell$, CM number field $L$, maximal real
subfield $L^+$, algebraic closure $\Omega$, and
$s:\mathrm{HidaSetting}(\ell,L,\Omega)$.  Put $S=s.S$.  The initial
level $V_0:s.\mathrm{Lev}$ satisfies $V_0.S'=S$.  Thus
\[
 \Gamma=G_{L^+,S}
  =G_{L^+}/N_S,\qquad G_0=\ker\eta,\qquad
 G'=\ker(\varepsilon|_{G_0}),
 \tag{CAP.1}
\]
where $N_S$ is the closed normal subgroup generated by all inertia
groups outside $S$.  The map $\rho:G_L\to\Gamma$ is the original
$V_0.\mathrm{pr}$ followed by the inclusion of $G_0$.
The fixed local datum is $\mathrm{locOf}(r)$; in particular its local
groups are the decomposition groups at the original valuations
$s.A(v)$ and $s.A(u)$, including their original inertia conditions.
Changing the global level below will leave this local datum fixed.

Let $O$ be any coefficient DVR satisfying the original
$\mathrm{IsCoef}(\ell,O)$, with its specified $\mathbb Z_\ell$-algebra
structure, and let $r:\mathrm{Resid}(V_0.\mathrm{frame},O,n)$ for
\emph{any} $n\geq0$.  Write $k=O/\mathfrak m_O$ and $\bar\chi_i$
for its residual characters.  The original hypothesis implies
$\operatorname{char}k=\ell$.  We use additive notation for
\[
 X_\ell(K)=\operatorname{Hom}_{\mathrm{cont}}(K,\mathbb F_\ell).
 \tag{CAP.2}
\]
Since the target is finite discrete, continuity is equivalent to an
open kernel.  Thus this is exactly the source's $\mathrm{HomL}$,
written additively.  The assertion $\Phi_\ell(\Gamma)$ means that
$X_\ell(K)$ is finite for \emph{every} open subgroup $K\leq\Gamma$.

The source anchors for the contracts are:
\begin{quote}\small
\nolinkurl{TaylorWilesObstructionBound.lean}:23--31,37--71,88--124;\\
\nolinkurl{ThetaCoefficientMapEq.lean}:12--40;\\
\nolinkurl{PairedDualFirstOrderDualMatrixSpace.lean}:78--106.
\end{quote}
All source paths and exact byte ranges are recorded in the accompanying
provenance; the proofs below establish the mathematical implications.

\subsection{Finiteness and the exact residual subgroup}

% BEGIN CAP1
\paragraph{CAP1. Finite first cohomology implies the full $\Phi$ condition.}
\label{claim:CAP1}
Let $\Pi$ be a profinite group.  Suppose $H^1_{\mathrm{cont}}(\Pi,M)$
is finite for every finite discrete $\mathbb F_\ell$-module $M$ with
continuous $\Pi$-action.  Then $\Phi_\ell(\Pi)$ holds.

\emph{Proof.}
Fix an open subgroup $K$.  The left-coset space $X=\Pi/K$ is finite:
its disjoint open cosets cover the compact group.  Choose representatives
$a_x$, with $a_K=1$, and let $M=\operatorname{Map}(X,\mathbb F_\ell)$,
with $(g m)(x)=m(g^{-1}x)$.  This is a finite continuous module.  For
$\chi\in X_\ell(K)$ define
\[
 z_\chi(g)(x)
   =\chi\bigl(a_x^{-1}g a_{g^{-1}x}\bigr).
 \tag{CAP.3}
\]
The argument belongs to $K$.  Inserting $a_{g^{-1}x}a_{g^{-1}x}^{-1}$
gives
\[
 z_\chi(gh)(x)=z_\chi(g)(x)+z_\chi(h)(g^{-1}x),
\]
so $z_\chi$ is a cocycle.  For each $x$, the representative
$a_{g^{-1}x}$ is locally constant in $g$; on each of its finitely many
clopen pieces (CAP.3) is a continuous map into $K$, followed by
$\chi$.  Hence $z_\chi$ is a continuous cocycle.

The map $\chi\mapsto[z_\chi]$ is injective.  Indeed, for $h\in K$,
$z_\chi(h)(K)=\chi(h)$, whereas every coboundary $h m-m$ has value
zero at the coset $K$.  If two cocycles define the same class, apply
this observation to their difference.  Thus $X_\ell(K)$ injects into
the finite group $H^1_{\mathrm{cont}}(\Pi,M)$.  This proves the result
for every open $K$, without requiring $K$ to be normal.\hfill$\square$

This is the written induced-cocycle argument underlying
\nolinkurl{ChebotarevFamilies/PhiGS.lean}:47--141 and
\nolinkurl{PhiDec.lean}:19--69; no separate Shapiro theorem is needed.
% END CAP1

% BEGIN CAP2
\paragraph{CAP2. The actual global and local $\Phi$ conditions.}
\label{claim:CAP2}
The original $\mathrm{CohP}$ supplies $\Phi_\ell(\Gamma)$ and
$\Phi_\ell(D_{s.A(v),L})$, $\Phi_\ell(D_{s.A(u),L})$ for all the
original local indices.

\emph{Proof.}
For the global assertion apply $\mathrm{CohP.gf}$ to the number field
$L^+$, the given $\Omega$, the finite set $V_0.S'$, and exponent $1$.
This set contains every prime above $\ell$, since $s.S\subseteq V_0.S'$.
The module in CAP1 is finite, discrete, and its action is continuous,
so it satisfies every hypothesis of
$\mathrm{GlobalFinitenessAt}$.  Its $H^1$ is therefore finite.  CAP1
applies to $G_{L^+,V_0.S'}$.

For the local assertion use $\mathrm{CohP.lf}$, again at exponent $1$,
with the number field $L$, the same $\Omega$, the actual valuation
$s.A(v)$, and its recorded prime $s.w(v)$.  The field extension and
place-over hypotheses are exactly those stored in $s$.  A decomposition
group is a closed subgroup of the absolute Galois group, hence
profinite.  CAP1 applies to it; the argument for $u\in s.St$ is
identical with its own recorded valuation.  These are the precise groups
in $\mathrm{locOf}(r)$, rather than newly chosen isomorphic models.
The other fields of $\mathrm{CohP}$ remain part of the original
hypothesis; none is replaced or asserted to follow here.\hfill$\square$

The exact finiteness predicates are
\nolinkurl{ReciprocityTheory/Hypotheses/GalCoh.lean}:17--20,151--155,
195--199,219--224.  The corresponding original applications are
\nolinkurl{PhiGS.lean}:142--160 and \nolinkurl{PhiDec.lean}:70--79.
% END CAP2

% BEGIN CAP3
\paragraph{CAP3. The original residual kernel is open and normal.}
\label{claim:CAP3}
There is an open normal subgroup $H\leq\Gamma$ with the exact formula
\[
 g\in H\quad\Longleftrightarrow\quad
 \eta(g)=1,\quad\bar\chi_i(g)=1\ (\forall i),\quad
 \bar\varepsilon(g)=1,
 \tag{CAP.4}
\]
where $\bar\varepsilon$ is reduction of $\varepsilon$ in
$\mathbb F_\ell^\times$.  In particular $H\subseteq G_0$ and both
the source $H_\eta$ and $H_\varepsilon$ clauses hold.

\emph{Proof.}
The displayed set is a subgroup.  It is open because $G_0$ is open,
the finitely many $\bar\chi_i$ have open kernels, and the residual
cyclotomic character is continuous with finite discrete target.
It is stable under conjugation by $G_0$, since all its conditions are
kernels of characters on $G_0$.  For conjugation by the specified
complex conjugation $c$, the residual determinant polarization is
essential.

Fix $x\in H$.  The image of $\varepsilon(x)$ in $k$ is $1$:
$\bar\varepsilon(x)=1$ says $\varepsilon(x)-1\in\ell\mathbb Z_\ell$,
which maps to zero in $k$.  The source involution consequently sends
$t[x]-1$ to $t[cxc^{-1}]^{-1}-1$.  Applying the full residual polarized
determinant law gives, for every $t\in k$,
\[
 (t-1)^n=
   \prod_{i=1}^{n}\bigl(t\bar\chi_i(cxc^{-1})^{-1}-1\bigr).
 \tag{CAP.5}
\]
For a particular index $j$ take $t=\bar\chi_j(cxc^{-1})$.
The right side is zero, so $(t-1)^n=0$.  Existence of $j$ implies
$n>0$; since $k$ is a field, $t=1$.  When $n=0$ there are no such
conditions.  The cyclotomic and $\eta$ conditions are preserved by
every conjugation.  Thus $cHc^{-1}=H$.  As
$\eta(\Gamma)=\{1,-1\}$ and $\eta(c)=-1$, every element of $\Gamma$
belongs to $G_0$ or $cG_0$, proving normality.

This retains exactly (CAP.4).  Taking an arbitrary normal core would
give an inclusion but would not establish the required equality
$\mathrm{TWPl.H\_eq}$.  No division by $n!$ has occurred; the argument
includes $\ell=3$ in rank four.\hfill$\square$

Source polarization:
\begin{quote}\small
\nolinkurl{UniversalOrdinaryDeterminantRing.lean}:80--99,109--117.\\
Residual-kernel consumer: \nolinkurl{ChebotarevFamilies/StarSingle.lean}:20--152.
\end{quote}
% END CAP3

\subsection{The precise relative pro-\texorpdfstring{$\ell$}{ell} subgroup and its tower}

For a subgroup $K\leq\Pi$ put
\[
 P_\ell^\Pi(K)=
 \bigcap_{\substack{U\triangleleft_{\mathrm{open}}\Pi\\
  \forall x\in K\ \exists a\geq0:\ x^{\ell^a}\in U}}U.
 \tag{CAP.6}
\]
This is the literal source $\mathrm{PL}(\ell,K)$.  Because $\Pi/U$
is finite, its displayed condition says precisely that the image of
$K$ is an $\ell$-group: every element has $\ell$-power order, and
Cauchy's theorem excludes every other prime divisor of its order.
Each term is closed and normal, so their intersection is closed and
normal.  The ambient group in the superscript matters.

% BEGIN CAP4
\paragraph{CAP4. Character kernels and local pro-$\ell$ kernels.}
\label{claim:CAP4}
If $K$ is open normal in a profinite group $\Pi$, every
$\chi\in X_\ell(K)$ vanishes on $P_\ell^\Pi(K)$.  Also, if
$f:I\to\Pi$ is continuous and $f(I)\subseteq K$, then
\[
 f\bigl(\operatorname{proKer}_\ell(I)\bigr)
       \subseteq P_\ell^\Pi(K),
 \tag{CAP.7}
\]
where $\operatorname{proKer}_\ell(I)$ is the intersection of all open
normal subgroups of $I$ of $\ell$-power index.

\emph{Proof.}
For the first assertion $A=\ker\chi$ is open in $\Pi$ and has finite
index.  Its normal core $U=\bigcap_{g\in\Pi}gAg^{-1}$ is an intersection
of finitely many conjugates, hence is open normal.  For $x\in K$ and
every $g\in\Pi$, normality of $K$ gives $g^{-1}xg\in K$, and
$\chi((g^{-1}xg)^\ell)=0$.  Hence $x^\ell\in U$.
Thus $U$ is one of the terms in (CAP.6), so
$P_\ell^\Pi(K)\subseteq U\subseteq\ker\chi$.

For (CAP.7), fix any term $U$ of (CAP.6).  The image of
$I\xrightarrow{f}\Pi\to\Pi/U$ is a finite $\ell$-group, as it lies
in the image of $K$.  Its kernel is open normal in $I$ and has
$\ell$-power index.  It therefore contains
$\operatorname{proKer}_\ell(I)$.  Intersecting over $U$ proves (CAP.7).
In the local application $I$ is the actual inertia group, and its
image $\operatorname{proKer}$ is exactly the source
$\mathrm{ineKer}$.\hfill$\square$

Compare \nolinkurl{ChebotarevFamilies/HomBound.lean}:5--78,
\nolinkurl{ProkerPL.lean}:5--32, and the definitions in
\texttt{GaloisCohomology\textsubscript{2}/ContinuousCycMod.lean}:10--36.
% END CAP4

% BEGIN CAP5
\paragraph{CAP5. A normal Frattini tower cofinal above the exact $P_\ell$.}
\label{claim:CAP5}
If $\Phi_\ell(\Gamma)$ and $H$ is open normal, there are descending
open normal subgroups $W_N\subseteq H$, $N\geq0$, such that
\[
 P_\ell^\Gamma(H)\subseteq V\triangleleft_{\mathrm{open}}\Gamma
 \quad\Longrightarrow\quad
 \exists N_0\ \forall N\geq N_0:\ W_N\subseteq V.
 \tag{CAP.8}
\]

\emph{Proof.}
For an open normal $K$ define
$\operatorname{fr}(K)=\bigcap_{\chi\in X_\ell(K)}\ker\chi$.
This is a finite intersection by $\Phi_\ell(\Gamma)$, and hence is
open.  Conjugation by $\Gamma$ permutes $X_\ell(K)$, so it is normal
in $\Gamma$.  Set $W_0=H$ and $W_{N+1}=\operatorname{fr}(W_N)$.

Fix $V$ as in (CAP.8).  The family of open normal subgroups in (CAP.6)
is closed under finite intersections: in the quotient by an
intersection, the image of $H$ embeds into a product of finite
$\ell$-groups.  Compactness shows that some such subgroup $U$ lies
in $V$.  Indeed, otherwise their intersections with the compact set
$\Gamma\setminus V$ would have the finite intersection property and
a common point, contradicting $P_\ell^\Gamma(H)\subseteq V$.
Consequently the image of $H$ in $\Gamma/V$ is a finite $\ell$-group.

Every nontrivial finite $\ell$-group $Q$ has a nonzero homomorphism to
the additive group $\mathbb F_\ell$.  To see this without a structural
theorem, the class equation shows that its center has order divisible
by $\ell$.  Choose a central subgroup of order $\ell$ by Cauchy's
theorem.  If the quotient by it is nontrivial, induction on $|Q|$
gives the desired homomorphism through that quotient; if it is trivial,
$Q$ itself is cyclic of order $\ell$.

Apply this observation to the image of $W_N$ in $\Gamma/V$.  Whenever
that image is nontrivial, a nonzero character of it pulls back to a
continuous character of $W_N$.  Hence the image of
$\operatorname{fr}(W_N)$ lies in a proper kernel and has strictly
smaller order.  After finitely many steps the image is trivial.
The descending property then gives (CAP.8) for every later $N$.
This is cofinality relative to $P_\ell^\Gamma(H)$, which is exactly
the condition used by $\mathrm{TWGrp.ph\_lim}$.\hfill$\square$

The source construction is \nolinkurl{ChebotarevFamilies/FrSNormal.lean}:6--207.
% END CAP5

\subsection{Chebotarev with the required places and sign}

% BEGIN CAP6
\paragraph{CAP6. A prescribed coset has degree-one arithmetic Frobenius.}
\label{claim:CAP6}
Let $E$ be a number field, $\Omega$ an algebraic closure, $U$ an open
subgroup of $G_E$, $g\in G_E$, and $A$ a finite set of finite primes of
$E$.  There are a prime $w\notin A$ of absolute norm a rational prime
$p$, a valuation $B$ of $\Omega$ above $w$, and
$\sigma\in D_{B,E}$ with
\[
 \sigma\text{ arithmetic Frobenius at }B,\qquad
 \sigma g^{-1}\in U.
 \tag{CAP.9}
\]

\emph{Published input and hypotheses.}
We use Chebotarev for a finite Galois extension of number fields:
each conjugacy class occurs as an unramified Frobenius class at a
positive density of primes.  The primary reference is J.~S.~Milne,
\emph{Algebraic Number Theory}, v3.08 (19 July 2020), Theorem 8.31,
printed p.~147, \url{https://www.jmilne.org/math/CourseNotes/ANT.pdf}.
Its Frobenius convention is the residue-field power map, as in
the discussion preceding 8.14.  We apply the theorem below to the
explicit finite Galois extension $E'/\mathbb Q$; no effective estimate
or RH hypothesis is involved.

\emph{Proof of the matching.}
Choose an open normal subgroup $U_0\subseteq U$; the normal core of
an open subgroup suffices.  Its fixed field $F/E$ is finite Galois.
Choose a finite Galois extension $E'/\mathbb Q$ inside $\Omega$
containing $F$, and put $h=g|_{E'}$.  Since $h$ fixes $E$, it belongs
to $\operatorname{Gal}(E'/E)$.  Chebotarev supplies infinitely many
rational primes unramified in $E'$ with Frobenius conjugate to $h$.
Exclude the finitely many rational primes below $A$.  Conjugating a
prime of $E'$ if necessary, choose $\mathfrak P$ whose Frobenius is
\emph{exactly} $h$.  Put $w=\mathfrak P\cap O_E$.
For every $a\in O_E$, the equality $h(a)=a$ implies
$a^p\equiv a\pmod w$.  Thus the residue field of $w$ is
$\mathbb F_p$ and $\operatorname{Norm}w=p$.

Extend $\mathfrak P$ to a valuation $B$ of $\Omega$ and lift the
arithmetic residue Frobenius to $D_{B,\mathbb Q}$.
For completeness, the finite valuation facts and this limiting passage
can be proved as follows.  In a finite Galois extension $F/K$, suppose
primes $P,Q$ above the same prime were in different Galois orbits.
The Chinese remainder theorem gives an integer $b$ that is zero
modulo $Q$ and one modulo every conjugate of $P$.  Its norm belongs
to $Q\cap O_K$ and hence to $P$.  Since the norm is a product of
conjugates of $b$, primality of $P$ contradicts those congruences.
Thus the action on primes is transitive.

The decomposition group at $P$ induces the full residue Galois
group.  To verify surjectivity explicitly, choose an integer $z$
whose residue at $P$ generates the multiplicative group of the finite
residue field, and whose residue at every other prime above the base
prime is zero.  Put $f(T)=\prod_{g\in\operatorname{Gal}(F/K)}(T-gz)$.
Its coefficients lie in $O_K$.  If the base residue cardinality is
$q$, then $f(z^q)=0$ modulo $P$.  Some $g$ therefore satisfies
$gz\equiv z^q\pmod P$.  Such a $g$ must stabilize $P$, since otherwise
$gz\equiv0\pmod P$.  Its residue action is the $q$-power map, which
generates the residue Galois group.  The kernel of this residue action
is exactly inertia by definition.

Finite integral lying-over and passage through finite composita give
a valuation $B$ extending $\mathfrak P$.  Fix this valuation.  For
each finite Galois subextension of $\Omega/\mathbb Q$ containing
$E'$, require an automorphism to stabilize its restriction of $B$
and induce the $p$-power residue map.  These are nonempty closed
conditions on the compact absolute Galois group by the preceding
finite argument.  Any finitely many are simultaneously satisfied
in their compositum.  Compactness supplies an automorphism satisfying
all of them; it stabilizes $B$ and induces the $p$-power map on every
residue.  The same compactness argument applied to finite prime
transitivity gives transitivity on valuations over a fixed prime.
Because $E'/\mathbb Q$ is unramified at $\mathfrak P$, this lift
$\sigma$ restricts to $h$ on $E'$.  It fixes $E$, belongs to
$D_{B,E}$, and induces the $\operatorname{Norm}w$ power map on
residues.  Finally $\sigma g^{-1}$ fixes $F$, so it belongs to
$U_0\subseteq U$.\hfill$\square$

The source target is \nolinkurl{ChebotarevFamilies/ChebOne.lean}:54--108;
the degree-one input is traced in
\nolinkurl{Proofs/InertiaEigenvalues/ChebotarevDegreeOne.lean}:17--69.
The argument above derives the required degree-one conclusion directly
from the stated finite Galois theorem.
% END CAP6

% BEGIN CAP7
\paragraph{CAP7. Split CM places and every prescribed finite list.}
\label{claim:CAP7}
In the original CM field $L$, fix an open subgroup $U\leq G_L$,
a finite set $A$ of primes of $L^+$, and a list
$g_j\in G_L$, $j\in\operatorname{Fin}(q)$.  One can choose
$v_j,w_j,B_j,\sigma_j$ such that the $v_j$ are distinct and outside
$A$, $B_j$ lies over $w_j$ and $v_j$,
\[
 D_{B_j,L^+}\subseteq G_L,
 \qquad \operatorname{Norm}w_j=p_j\text{ is prime},\qquad
 \sigma_jg_j^{-1}\in U,
 \tag{CAP.10}
\]
and $\sigma_j$ is an arithmetic Frobenius at $B_j$ over $L$.

\emph{Proof.}
Choose an algebraic integer $a\in O_L$ with $a\ne\bar a$; a
nonreal element of $L$ becomes integral after multiplication by a
nonzero rational integer.  Enlarge the finite forbidden set of primes
of $L$ to include every prime above $A$ and every prime dividing
$a-\bar a$.  Finiteness follows from finite extension degree and the
factorization of the nonzero principal ideal $(a-\bar a)$.
Apply CAP6.  If $B$ is the resulting valuation and $w$ its prime,
an element $d\in D_{B,L^+}$ restricts on $L$ to either $1$ or CM
conjugation.  It stabilizes $w$ and acts trivially on its residue field
$\mathbb F_p$.  The second possibility would imply
$a-\bar a\in w$, which was excluded.  Thus it fixes $L$, proving
the decomposition-group inclusion.  The prime $v$ below $w$ is
outside $A$.

For the list choose its entries in order, adding previously chosen
$v_i$ to $A$ at each step.  This gives injective $j\mapsto v_j$ and
retains all coset conditions.  When $q=0$ use the empty family.

We will also use the following consequence of the valuation facts
in CAP6.  If $B$ and $B'$ lie over the same prime of $L^+$, they are
conjugate by $G_{L^+}$: the finite-level transitivity statements have
a simultaneous solution by compactness.  The subgroup $G_L$ is
normal in $G_{L^+}$ because $L/L^+$ is quadratic Galois.  Therefore
$D_{B,L^+}\subseteq G_L$ implies
$D_{B',L^+}\subseteq G_L$.  This supplies splitting at the original
stored valuation $s.A(v)$ even when the prime was selected using a
different valuation $B$.\hfill$\square$

These are the exact conditions of \nolinkurl{ChebOne.lean}:6--108 and
\nolinkurl{ChebFam.lean}:6--35.
% END CAP7

\subsection{The actual level maps and the uniform bound}

% BEGIN CAP8
\paragraph{CAP8. Enlarging the original level with its fixed local datum.}
\label{claim:CAP8}
Let $Q$ be a finite set of primes of $L^+$, each with a valuation
$B_v$ over $v$ satisfying $D_{B_v,L^+}\subseteq G_L$.  There is an
actual level $V_1:s.\mathrm{Lev}$ with
$S_1=V_1.S'=V_0.S'\cup Q$ and a surjective continuous map
\[
 p:G_{L^+,S_1}\longrightarrow G_{L^+,V_0.S'},\qquad
 p(\operatorname{toGS}_{S_1}g)=\operatorname{toGS}_{V_0.S'}g,
 \tag{CAP.11}
\]
which is a $\mathrm{Glob.Hom}$ from
$\mathrm{globOf}(r,V_1)$ to $\mathrm{globOf}(r,V_0)$.

\emph{Proof.}
Since $V_0.S'\subseteq S_1$, the generators of $N_{S_1}$ are among
those of $N_{V_0.S'}$, so $N_{S_1}\subseteq N_{V_0.S'}$.
The quotient map gives (CAP.11), including surjectivity and continuity.
Set $\eta_1=\eta_0p$ and $\varepsilon_1=\varepsilon_0p$.
The original $\eta$ characterization, value at $c$, open kernel,
cyclotomic identity, and continuity pull back along $p$.
The enlarged set is finite and contains $s.S$.  At old primes the
stored splitting condition is unchanged.  At new primes CAP7
transports the supplied splitting condition to $s.A(v)$.
These verify every field of $V_1$.

The global datum uses the same $\mathrm{locOf}(r)$ on both sides.
Its local maps are restrictions of the quotient maps followed by
$G_L\to G_{L^+}$.  Consequently (CAP.11) commutes with every
original $\varphi_v$ and $\varphi_u$, as well as $\eta$, $c$, and
$\varepsilon$.  These are precisely all fields of
$\mathrm{Glob.Hom}$.  In particular its induced map on $G_0$
satisfies $p_0(V_1.\mathrm{pr}(g))=V_0.\mathrm{pr}(g)$ for
every $g\in G_L$.\hfill$\square$

Source level construction and contracts:
\begin{quote}\small
\nolinkurl{ChebotarevFamilies/LevUp.lean}:5--53;\\
\nolinkurl{ContinuousResP.lean}:106--162;\\
\nolinkurl{PairedDualFirstOrderDualMatrixSpace.lean}:78--99.
\end{quote}
% END CAP8

% BEGIN CAP9
\paragraph{CAP9. The whole level kernel is generated modulo the exact $P_\ell$.}
\label{claim:CAP9}
Suppose $Q=\{v_j:j\in\operatorname{Fin}(q)\}$ is disjoint from
$V_0.S'$, with the valuations $B_j$ of CAP7.  Let $p$ be (CAP.11),
$K=p^{-1}H$, and let $\tau_j\in G_L$ satisfy
$\mathrm{IsProGen}(L,B_j,\ell,\tau_j)$.  Put
$t_j=V_1.\mathrm{pr}(\tau_j)$, regarded also as an element of
$\Gamma_1=G_{L^+,S_1}$.  Then $p(t_j)=1$ and
\[
 \ker p\ \subseteq\
 \overline{\langle\!\langle t_j:j\in\operatorname{Fin}(q)\rangle\!\rangle
       \ P_\ell^{\Gamma_1}(K)}.
 \tag{CAP.12}
\]
The bar means topological closure of the subgroup generated by the two
displayed normal subgroups.  Moreover the image of every
$\mathrm{ineKer}(L,B_j,\ell)$ lies in $P_\ell^{\Gamma_1}(K)$.

\emph{Proof.}
Every inertia element at $v_j$ is killed by the old quotient, because
$v_j\notin V_0.S'$.  Hence its image in $\Gamma_1$ is in $\ker p$
and therefore in $K$.  Apply CAP4 to the continuous map from its
actual inertia group to $\Gamma_1$: its pro-$\ell$ kernel maps into
$P_\ell^{\Gamma_1}(K)$.  The definition of $\mathrm{IsProGen}$ says
that the inertia group is contained in the closure of the subgroup
generated by $\tau_j$ and its $\mathrm{ineKer}$.  Its whole image
therefore lies in the right side $C$ of (CAP.12); continuity suffices
because $C$ is closed.  In particular $p(t_j)=1$.

The group $C$ is closed normal in $\Gamma_1$.  To show that it kills
the complete old ramification kernel, take an inertia group at a
prime $v\notin V_0.S'$, with any valuation $A$ above $v$.
If $v\notin S_1$ it is killed already by $\Gamma_1$.
Otherwise $v=v_j$.  By valuation transitivity choose a conjugation
carrying $A$ to $B_j$.  At $B_j$ the decomposition group over $L^+$
is contained in $G_L$, so its inertia is exactly the inertia over
$L$, with the same residue action.  The previous paragraph puts its
image in $C$, and normality transports this conclusion back to $A$.
Thus the inverse image of $C$ in $G_{L^+}$ contains every inertia
generator of $N_{V_0.S'}$.  It is closed normal, so contains all of
$N_{V_0.S'}$.  Every element of $\ker p$ lifts to an element of this
last group, which proves (CAP.12).\hfill$\square$

This accounts for the full, closed kernel in
\nolinkurl{ChebotarevFamilies/PLNormal.lean}:37--101, including arbitrary
valuations over each new prime.
% END CAP9

% BEGIN CAP10
\paragraph{CAP10. A character bound uniform in the enlarged level.}
\label{claim:CAP10}
Let $p:\Pi\twoheadrightarrow\Gamma$ be continuous and surjective
between profinite groups, let $H\triangleleft_{\mathrm{open}}\Gamma$,
and set $K=p^{-1}H$.  Suppose $X_\ell(H)$ is finite and $t_1,\ldots,t_q$
lie in $\ker p$, with (CAP.12) for this $p$ and $K$.  Then
\[
 |X_\ell(K)|\leq |X_\ell(H)|\,
       \ell^{\,q[\Gamma:H]}.
 \tag{CAP.13}
\]
In particular $X_\ell(K)$ is finite.

\emph{Proof.}
The quotient map identifies $\Pi/K$ with $\Gamma/H$.  Choose
representatives $a_1,\ldots,a_d$ for these cosets, where
$d=[\Gamma:H]$.  Evaluate a character on the $qd$ elements
$a_i t_j a_i^{-1}\in K$.  There are at most $\ell^{qd}$ possible
arrays of values.

If two characters have the same array, their difference $\chi$
vanishes on these elements and hence on every conjugate of every
$t_j$.  Indeed, write any $a\in\Pi$ as $k a_i$ with $k\in K$;
conjugation by $k$ does not affect an abelian character of $K$.
It vanishes on $P_\ell^\Pi(K)$ by CAP4.  Its kernel is closed in
$\Pi$, since $K$ is open and closed and the character is continuous.
The kernel therefore contains the right side of (CAP.12), so
$\chi$ kills $\ker p$.

It descends uniquely to a character $\chi_0:H\to\mathbb F_\ell$.
The restricted surjection $p:K\to H$ is a quotient map: a continuous
surjection from a compact space onto a Hausdorff space has this
property.  Consequently $\chi_0$ is continuous, since its composite
with $p$ is $\chi$.  In every nonempty evaluation fiber choose one
character; subtracting it injects that fiber into $X_\ell(H)$.
There are at most $\ell^{qd}$ fibers, proving both finiteness and
(CAP.13).  The proof includes $q=0$.

Only $q,H,\ell$ occur in the bound.  In particular one may use the
same integer
$h_b=|X_\ell(H)|\ell^{q[\Gamma:H]}$ at every level.\hfill$\square$

This proves the actual conclusion of
\nolinkurl{ChebotarevFamilies/HomBound.lean}:82--157.
% END CAP10

\subsection{Construction of every field of the original records}

% BEGIN CAP11
\paragraph{CAP11. Simultaneous congruence, approximation, and inverse Frobenius.}
\label{claim:CAP11}
Fix an arbitrary list $s_j\in G'$, $j\in\operatorname{Fin}(q)$.
There are, for every $N\geq0$, distinct new split primes $v_{Nj}$,
primes $w_{Nj}$ of $L$, valuations $B_{Nj}$, arithmetic Frobenius
elements $\sigma_{Nj}$, and tame generators $\tau_{Nj}$, such that
\[
 \begin{gathered}
 v_{Nj}\notin S,\qquad
 Q_{Nj}:=\operatorname{Norm}w_{Nj}\equiv1\pmod{\ell^N},\\
 \rho(\sigma_{Nj})^{-1}s_j^{-1}\in W_N.
 \end{gathered}
 \tag{CAP.14}
\]
With the actual level $V_N$ attached to
$S_N=S\cup\{v_{Nj}\}_j$, write $p_N$ for its map to the original
level, $K_N=p_N^{-1}H$, and
\[
 \phi_{Nj}=V_N.\mathrm{pr}(\sigma_{Nj})^{-1},\qquad
 t_{Nj}=V_N.\mathrm{pr}(\tau_{Nj}).
 \tag{CAP.15}
\]
Then $p_N(t_{Nj})=1$, (CAP.12) holds, and
\[
 (t_{Nj}^{Q_{Nj}})^{-1}
    (\phi_{Nj}^{-1}t_{Nj}\phi_{Nj})
       \in P_\ell^{\Gamma_N}(K_N).
 \tag{CAP.16}
\]

\emph{Proof.}
First, the original $V_0.\mathrm{pr}:G_L\to G_0$ is surjective.
Lift $s_j\in G_0$ through the quotient $G_{L^+}\to\Gamma$; its
$\eta$ value is $1$, so the exact $V_0.h\eta$ identifies the lift
as an element of the image of $G_L$.  Choose $g_j\in G_L$ with
$V_0.\mathrm{pr}(g_j)=s_j$.  Since $s_j\in G'$ and the original
$V_0.h\varepsilon$ is the full cyclotomic identity,
$\varepsilon_\ell(g_j)=1$ in $\mathbb Z_\ell^\times$, not merely
after residual reduction.

Set
\[
 U_N=\rho^{-1}(W_N)\cap\ker(\varepsilon_\ell\bmod\ell^N).
 \tag{CAP.17}
\]
This is open: $W_N$ is open, and the second character is continuous
with finite discrete target.  Apply CAP7 to $U_N$, the forbidden
set $S$, and the list $g_j^{-1}$.  The result satisfies
$\sigma_{Nj}g_j\in U_N$.  In particular
$\rho(\sigma_{Nj})s_j\in W_N$.  Normality of $W_N$, followed by
inversion, gives
\[
 \rho(\sigma_{Nj})^{-1}s_j^{-1}
  =\bigl(s_j(\rho(\sigma_{Nj})s_j)s_j^{-1}\bigr)^{-1}\in W_N,
\]
with the precise order of factors required in (CAP.14).

The prime $w_{Nj}$ is away from $\ell$, because otherwise its prime
$v_{Nj}$ below it would belong to $S$.  The cyclotomic image of an
arithmetic Frobenius there is $Q_{Nj}$ modulo $\ell^N$.
Here is the normalization check.  An $\ell^N$-th root of unity has
valuation one.  Arithmetic Frobenius sends its residue to its
$Q_{Nj}$-th power.  Reduction is injective on these roots: if
$u^{\ell^N}=1$ and $u\equiv1$ in the residue field, then
$(u-1)(1+u+\cdots+u^{\ell^N-1})=0$, and the second factor reduces
to the unit $\ell^N$, forcing $u=1$.  Thus Frobenius acts on the
roots themselves by the $Q_{Nj}$-th power, which is the claimed
cyclotomic value.  On the other hand
$\varepsilon_\ell(g_j)=1$ and $\sigma_{Nj}g_j\in U_N$ imply that
this value is $1$.  This proves the congruence, including the trivial
case $N=0$.

Now apply the retained $\mathrm{CohP.tame.1}$ to the number field
$L$, the given algebraic closure, $B_{Nj}$ over $w_{Nj}$, and the
prime $\ell$ which does not lie in $w_{Nj}$.  Its
$\mathrm{TameQuotient}$ clause supplies
$\tau_{Nj}$ with $\mathrm{IsProGen}$ and the exact relation
\[
 (\tau_{Nj}^{Q_{Nj}})^{-1}
   \sigma_{Nj}\tau_{Nj}\sigma_{Nj}^{-1}
       \in\mathrm{ineKer}(L,B_{Nj},\ell).
 \tag{CAP.18}
\]
CAP8 supplies $V_N$ and $p_N$.  CAP9 supplies its complete kernel
bound, $p_N(t_{Nj})=1$, and the inclusion of $\mathrm{ineKer}$ into
$P_\ell^{\Gamma_N}(K_N)$.  Mapping (CAP.18) to that quotient and
using $\phi_{Nj}^{-1}=V_N.\mathrm{pr}(\sigma_{Nj})$ gives exactly
(CAP.16).  The hypothesis $\mathrm{CohP}$ was retained throughout;
the tame relation is an explicit use of its original clause.\hfill$\square$

Source assembly and normalization:
\begin{quote}\small
\nolinkurl{AutomorphicDetection/ArPl.lean}:21--119;\\
\nolinkurl{ReciprocityTheory/GaloisValuations/Global.lean}:199--229;\\
tame contract: \nolinkurl{ContinuousCycMod.lean}:116--134.
\end{quote}
% END CAP11

% BEGIN CAP12
\paragraph{CAP12. The complete original $\mathrm{TWGrp}$ at all levels.}
\label{claim:CAP12}
For the fixed arbitrary list of CAP11, define $T$ by taking its
length to be $q$, its subgroup to be the exact $H$ of CAP3, its
listed elements to be $s_j$, and its level datum to be
$G_N=\mathrm{globOf}(r,V_N)$, with $h_N=p_N$ and the data
(CAP.14)--(CAP.15).  Give it the single bound
\[
 h_b=|X_\ell(H)|\ell^{q[\Gamma:H]}.
 \tag{CAP.19}
\]
These data satisfy every field of the original $\mathrm{TWGrp}$.

\emph{Proof.}
CAP3 provides $H$ open normal, $H_\eta$, and $H_\varepsilon$.
CAP8 provides every required $\mathrm{Glob.Hom}$ $h_N$ with the
same original local datum.  The list belongs to $G'$ by assumption.
The norm congruence is CAP11.  Its last relation in (CAP.14), together
with $W_N\subseteq H$, gives $\mathrm{ph\_H}$ at every $N$.
For any open normal $V\supseteq P_\ell^\Gamma(H)$ and any $j$,
CAP5 gives $W_N\subseteq V$ for all sufficiently large $N$; hence
$h_N(\phi_{Nj})s_j^{-1}\in V$ eventually.  This is the exact
$\mathrm{ph\_lim}$ quantifier, with $V$ and $j$ fixed before the
eventual bound.  In fact the constructed bound works for all $j$
at once.

The conditions $\mathrm{ta\_ker}$, $\mathrm{tame}$, and
$\mathrm{ker\_le}$ are CAP9 and CAP11.  Finiteness of $X_\ell(H)$
follows from CAP2.  CAP10 applied to every $p_N$ gives finiteness
of $X_\ell(K_N)$ and its bound by the same (CAP.19), which proves
both $\mathrm{fin0}$ and $\mathrm{finN}$.  Finally CAP2 supplies
$\mathrm{phiV}$ and $\mathrm{phiU}$ for all the original local
decomposition groups.  This exhausts the record fields in
\nolinkurl{ThetaCoefficientMapEq.lean}:15--40.

All $N\in\mathbb N$ are included, with primes and generators
chosen for each $N$.  The required compatibility is the map from
each level to the same original level, together with the displayed
uniform bound and eventual approximation.  No nestedness of the
sets of new primes is used.\hfill$\square$
% END CAP12

% BEGIN CAP13
\paragraph{CAP13. Actual arithmetic realization and the full original $\mathrm{ArPl}$.}
\label{claim:CAP13}
The system $T$ of CAP12 has an actual $\mathrm{TWPl}(r,T)$,
with its original length and list unchanged.  Consequently the entire
source predicate $\mathrm{ArPl}$ holds.

\emph{Proof.}
For every $N$ take the arithmetic level to be the already constructed
$V_N$.  Because $T.G_N=\mathrm{globOf}(r,V_N)$, take its comparison
$e_N$ to be the identity $\mathrm{Glob.Hom}$.  It is injective, and
its compatibility with $T.h_N$ is exactly (CAP.11), for every
$g\in G_{L^+}$.
Use the same $v_{Nj},w_{Nj},B_{Nj}$ of CAP11.  Their injectivity,
avoidance of $S$, actual place-over relations, and decomposition-group
splitting were proved in CAP7.  By construction
$V_N.S'=S\cup\operatorname{range}(v_N)$ and
$T.Nq_{Nj}=\operatorname{Norm}w_{Nj}$.

Set the record's Frobenius element to
$\mathrm{fr}_{Nj}=\sigma_{Nj}^{-1}$.  It belongs to the decomposition
group, and its \emph{inverse} is the arithmetic Frobenius.
This is precisely $\mathrm{fr\_frob}$ in the source.  Use the same
$\tau_{Nj}$ from CAP11, with its proved $\mathrm{IsProGen}$.
The comparisons for $\mathrm{ph}$ and $\mathrm{ta}$ are the
literal identities (CAP.15), since $e_N$ is the identity.
Finally $\mathrm{H\_eq}$ is (CAP.4), using the original residual
characters and the original cyclotomic character.  These checks cover
every field of \nolinkurl{TaylorWilesObstructionBound.lean}:88--111.

Before choosing the list, CAP2 already proved $\Phi_\ell(\Gamma)$.
For every $q$ and every list in $G'$, CAP12 and this construction now
give $T$ and its nonempty arithmetic realization.  The length equality
$T.q=q$ is definitional and its transport of the listed elements is
the identity.  This is exactly the conjunction and quantifier order
in the original \texttt{ArPl} definition at lines 119--124.
The argument permits $n=0$, $q=0$, and $N=0$; in the CountII
application it therefore applies in particular to $n=4$, $\ell=3$.
Besides the explicitly retained coefficient, level and residual data,
its original arithmetic hypothesis is the unchanged
$\mathrm{CohP}$.\hfill$\square$
% END CAP13

% BEGIN CAP14
\paragraph{CAP14. Contribution to the same-ideal CountII producer.}
\label{claim:CAP14}
The complete residual input \texttt{JSP\_ACTUAL\_ARPL} in CIP12 is
supplied by CAP1--CAP13.  The remaining sufficient same-ideal support
rule is JSP11 with CIP10--CIP11 and CAP13 inserted, retaining the
original \texttt{Choice}, first-order bound, obstruction bound,
\texttt{Raw}, and same-$J$ \texttt{HasAut}, and all its original
\texttt{CohP}/\texttt{LocCFT}/coefficient/level/cyclotomic hypotheses.

\emph{Proof and scope.}
The original consumer first uses the $\Phi$ conjunct of $\mathrm{ArPl}$
to call $\mathrm{Choice}$; after Choice produces its finite list, it
uses the second conjunct to realize precisely that list at all levels.
CAP2 is independent of the list, and CAP13 is universal in it, so
this substitution introduces no circular dependence and no additional
restriction on the detector chosen by Choice.

ArPl has no ideal argument.  Therefore it works for every already
fixed original $B,r,J$, without replacing $J$ when a point, list, or
level is chosen.  The same-$J$ HasAut premise is still the original
universal assertion for every realized $\mathrm{TWGrp}$ and
$\mathrm{TWPl}$; constructing these records does not construct their
automorphic modules.  The original ordinary/control and separating
arithmetic-fiber data, the four further Hida fields, and this HasAut
producer remain separate obligations.

CAP13 is universal over the original coefficient ring, residual datum,
and Hida setting.  It adds no coefficient enlargement and no threshold.
Thus it can be inserted for each of both original rank-four Links
after the fixed $\iota$ and their one common coefficient threshold,
for every later coefficient/Space/universal datum.  Within each Link
the ideal stays the original $J$ of JSP1--JSP12; the two different
universal rings are not identified.  The common $\kappa$, the actual
distinct auxiliary primes required by LinkTwo, and $\ell=3$ are
unchanged.  QX4/QX10 and YX2 therefore retain exactly their original
conditional support-to-CountII application.  CountII remains needed
for RH and unrestricted Dirichlet GRH.  KA5--KA6 remove only CaseIII
for targets of two-power order; Cases I, II, and V retain their
arithmetic producers.\hfill$\square$

The exact consumer order is
\nolinkurl{ChebotarevFamilies/PFromAutOf.lean}:30--49,104--126;
the joint producer contracts are
\nolinkurl{DenseThresholdCases/HidaU.lean}:12--16,
\nolinkurl{PointTheoremBaseObject.lean}:26--48,
\nolinkurl{LinkTwo.lean}:13--51, and
\nolinkurl{PotentialModularityTower.lean}:38--66.
% END CAP14

\subsection{The complete residual Hecke-value subpacket}

For this subsection retain an arbitrary original level $V:s.\mathrm{Lev}$
and $Q:\mathrm{QData}(V)$, with its original finite auxiliary set.
Here no equality $V.S'=s.S$ is imposed.  Fix $n\geq1$, a coefficient
field datum $c$, and \emph{any supplied} representation
$\tau_0:K_T(n,\mathrm{wt}(s))\to\mathrm{GL}_m(c.O)$.
Let $r:\mathrm{Resid}(V.\mathrm{frame},c.O,n)$ and put
$k=c.O/\mathfrak m$, a finite field of characteristic $\ell$.
The original $\mathcal H=\mathrm{hset}(Q,\tau_0)$ has operator indices
\[
 \mathcal H.\mathrm{Op}
   =(\mathcal H.W_p\times\operatorname{Fin}(n))
      \sqcup(V.\mathrm{frame}.V_\ell\times\operatorname{Fin}(n)),
 \quad
 \mathcal H.W_p=
 \{w:w|_{L^+}\notin V.S',\ w\text{ split}\}.
 \tag{CAP.20}
\]
We construct one function $a:\mathcal H.\mathrm{Op}\to c.O$ with
both original fields of $\mathrm{IsRes}(Q,r,\tau_0,a)$.
The validity of the \emph{actual} $\mathrm{artOf}$ on each original
chosen local group is the already proved CIP4.

% BEGIN CAP15
\paragraph{CAP15. Every residual local character is constant on Artin fibers.}
\label{claim:CAP15}
Let $D$ be the decomposition group at a number-field prime, $W$ its
original Weil subgroup, and $b:W\twoheadrightarrow K^\times$ any map
satisfying the full original $\mathrm{IsLocalArtin}$.  If
$\psi:D\to M$ is an abelian-valued character with open kernel, then
\[
 b(x)=b(y)\quad\Longrightarrow\quad\psi(x)=\psi(y)
 \qquad(x,y\in W).
 \tag{CAP.21}
\]
This applies to all the original $r.\bar\psi_{v,i}$ and to the
residual character $\mathcal H.\varepsilon_v$.

\emph{Proof.}
The image $B=\psi(D)$ is finite, because its kernel is open in the
compact group $D$.  Characters $B\to\mathbb C^\times$ separate its
elements.  An elementary proof of the latter assertion is useful:
for $1\ne z\in B$, define a character on $\langle z\rangle$ by
sending $z$ to a primitive root of its order.  Extend a character
from a subgroup $C$ to $\langle C,t\rangle$ as follows.  If $e$ is
the least positive integer with $t^e\in C$, choose
$u\in\mathbb C^\times$ with $u^e=\chi(t^e)$ and set
$\chi'(ct^j)=\chi(c)u^j$.  The minimality of $e$ and this equation
prove independence of the expression and multiplicativity.  Iterating
over the finite group extends the original character to $B$.

For any such complex character $\lambda$, the composite
$\lambda\psi|_W$ is smooth: its kernel on inertia contains the
open kernel of $\psi$ restricted to inertia.  The literal
$\mathrm{IsLocalArtin.factor}$ therefore writes it as $\xi\circ b$
for a smooth character $\xi$ of $K^\times$.  If $b(x)=b(y)$, it
follows that every $\lambda$ is $1$ on $\psi(xy^{-1})$.
Character separation gives $\psi(x)=\psi(y)$.

The kernels of $r.\bar\psi_{v,i}$ are open by the original Resid
record.  For $\mathcal H.\varepsilon_v$, the continuous cyclotomic
character modulo $\ell$ has open kernel, and that kernel maps to $1$
in $k^\times$: the specified map $\mathbb Z_\ell\to k$ kills
$\ell\mathbb Z_\ell$.  Hence this residual character also has open
kernel.  CIP4 supplies every required field of the actual
$b=\mathrm{artOf}(s.hw(v))$, so (CAP.21) applies without replacing it
by a preferred reciprocity map.\hfill$\square$

The original local arguments are \nolinkurl{DenseThresholdCases/TResEx.lean}:65--104;
the factorization predicate is
\texttt{AutomorphicForms\textsubscript{2}/HasExponentZeroIff.lean}:145--152.
% END CAP15

% BEGIN CAP16
\paragraph{CAP16. The exact residual Hecke polynomial, independent of every Frobenius choice.}
\label{claim:CAP16}
For each $w\in\mathcal H.W_p$ there are elements
$b_{w,i}\in c.O$, $i\in\operatorname{Fin}(n)$, such that for
\emph{every} valuation $B$ over $w$ and every arithmetic Frobenius
$\sigma$ at $B$, the original residual polynomial satisfies
\[
 \bar D_r\bigl(X-[V.\mathrm{pr}(\sigma^{-1})]\bigr)
   =\sum_{j=0}^{n}
       (-1)^j q_w^{j(j-1)/2}\,\beta_{w,j} X^{n-j},
 \qquad
 \beta_{w,0}=1,\quad \beta_{w,j}=\bar b_{w,j-1}\ (1\leq j\leq n),
 \tag{CAP.22}
\]
where $q_w=\operatorname{Norm}w$.  The left side denotes the exact
source $\bar D_r.\chi([V.\mathrm{pr}(\sigma^{-1})])$.

\emph{Proof.}
Every character $\bar\chi_i\circ V.\mathrm{pr}$ is unramified at
$w$: since the prime below $w$ is outside $V.S'$, the original
$N_{V.S'}$ kills inertia at \emph{every} valuation over $w$.
For two valuations $B,B'$ over $w$, choose $g\in G_L$ carrying
$B$ to $B'$ by valuation transitivity.  Conjugate a Frobenius at
$B'$ by $g^{-1}$ to get a Frobenius at $B$.  Its ratio with a fixed
Frobenius there is an inertia element, as both have the same residue
power action.  An unramified abelian character kills this ratio and
is invariant under conjugation.  Thus the value of each
$\bar\chi_i(V.\mathrm{pr}(\sigma^{-1}))$ is independent of both
$B$ and $\sigma$, with the inverse exactly as displayed.

Choose one such Frobenius $\sigma_w$ and form the monic degree-$n$
polynomial
\[
 P_w(X)=\prod_{i\in\operatorname{Fin}(n)}
     \bigl(X-\bar\chi_i(V.\mathrm{pr}(\sigma_w^{-1}))\bigr).
\]
It is the required determinant polynomial because the residual
representation is diagonal.  The prime $w$ is away from $\ell$:
every prime of $L^+$ above $\ell$ belongs to $s.S\subseteq V.S'$.
Thus $q_w$ is an $\ell$-adic unit and is nonzero in $k$.
For $1\leq j\leq n$ define
\[
 \beta_{w,j}=(-1)^j q_w^{-j(j-1)/2}
                         [X^{n-j}]P_w(X).
 \tag{CAP.23}
\]
Choose $b_{w,j-1}\in c.O$ lifting this element by surjectivity of
the residue map.  Substitution into the right side of (CAP.22)
recovers every coefficient of $P_w$: the leading coefficient is $1$,
and the two sign factors and two powers of $q_w$ cancel for every
$j\geq1$.  The Frobenius independence already proved gives (CAP.22)
for all choices.  This is the literal $\mathrm{heckePol}$
normalization, including its factor $q_w^{j(j-1)/2}$ and its
operator-index shift $j\mapsto j-1$.  No factorial or inverse of
$\ell$ is used.\hfill$\square$

The polynomial is defined in \nolinkurl{HeckeAlgebras/Distinct.lean}:46--47;
the corresponding source construction and independence arguments are
\nolinkurl{DenseThresholdCases/TResEx.lean}:16--64,109--140.
% END CAP16

% BEGIN CAP17
\paragraph{CAP17. Ordered uniformizer values and the entire original $\mathrm{IsRes}$.}
\label{claim:CAP17}
For every original $V,Q,c,m,\tau_0,r$ in the scope of (CAP.20), there
exists one $a:\mathcal H.\mathrm{Op}\to c.O$ satisfying the full
original $\mathrm{IsRes}(Q,r,\tau_0,a)$.

\emph{Proof.}
At each original $v\mid\ell$, CIP4 gives surjectivity of the
actual $\mathrm{artOf}$ onto $L_{s.w(v)}^\times$.  Choose a Weil
element $y_v$ whose image is the original chosen uniformizer and
regard it as an element of the original decomposition group.
For $j\in\operatorname{Fin}(n)$ set
\[
 u_{v,j}=\prod_{\substack{i\in\operatorname{Fin}(n)\\i\leq j}}
       r.\bar\psi_{v,i}(y_v)\,
               \mathcal H.\varepsilon_v(y_v)^i\ \in k^\times.
 \tag{CAP.24}
\]
CAP15 shows that every factor is unchanged if $y_v$ is replaced by
\emph{any} Weil element having that same Artin image.  Hence (CAP.24)
is independent of this choice.  Choose $d_{v,j}\in c.O$ whose residue
is the underlying field element of $u_{v,j}$.

Define the single function on the original disjoint union of operator
indices by
\[
 a(\mathrm{inl}(w,i))=b_{w,i},\qquad
 a(\mathrm{inr}(v,j))=d_{v,j}.
 \tag{CAP.25}
\]
For the first field of $\mathrm{IsRes}$, CAP16 proves precisely its
polynomial equation for every $w,B,\sigma$.  For the second, let
$x$ be any Weil element and $g$ any element of the recorded
decomposition group having the same underlying Galois automorphism,
with $\mathrm{artOf}(x)$ equal to the chosen uniformizer.  Then $g$
is the image of $x$ under the original Weil inclusion.  CAP15 replaces
every factor at $g$ by the corresponding factor at $y_v$.  Equation
(CAP.24), followed by the choice of $d_{v,j}$, is exactly the
source's ordered product over $i\leq j$, with exponent $i$, and
proves its \texttt{unif} field.  These are both fields of
\nolinkurl{HeckeAlgebras/IsTheta.lean}:161--168.

The choices of lifts are made simultaneously in the same original
ring $c.O$; no extension of coefficients is needed.  This proves
the complete conclusion of the original \texttt{TResEx} supplier,
at every original level and for every supplied $\tau_0$.  Its source
use of a literature hypothesis merely to obtain $\mathrm{artOf}$
is replaced here by CIP4 with its checked class field theory input.
In constructing $R:\mathrm{UnitaryHeckeData}$ this supplies its
literal fields $a$ and $\mathrm{res}$ after $\tau_0$ is supplied.
The chosen $a$ is then held fixed for the ordinary module and its
future kernel ideal.  No equality of modules or ideals associated
with different choices of $a$ is asserted.

The integral type with $\mathrm{TypeForAlt}$, the torus character
with both $\mathrm{IsTheta}$ and the full $\mathrm{IsWt}$ (including
its actual neat-level condition), and the remaining original
ordinary control, determinant, and arithmetic-fiber producers remain open at their
previous scopes.  In particular CAP17 does not supply all of $R$
or close $\mathrm{COUNT\_II\_SUPPORT}$.\hfill$\square$

The source conclusion is \nolinkurl{DenseThresholdCases/TResEx.lean}:105--165;
its precise insertion is \nolinkurl{UnitaryHeckeData.lean}:59--70 and
\nolinkurl{HdEx.lean}:20--27.  The retained torus-character and full
weight supplier is \nolinkurl{TThetaEx.lean}:160--180.
% END CAP17
